\documentclass{article}

% NeurIPS 2026 style
\usepackage[nonatbib]{neurips_2026}  % NeurIPS 2026 style — submission (blind review)

\usepackage[utf8]{inputenc}
\usepackage[T1]{fontenc}
\usepackage{hyperref}
\usepackage{url}
\usepackage{booktabs}
\usepackage{amsfonts}
\usepackage{amsmath}
\usepackage{amssymb}
\usepackage{nicefrac}
\usepackage{microtype}
\usepackage{graphicx}
\usepackage{xcolor}
\usepackage{subcaption}
\usepackage{algorithm}
\usepackage{algorithmic}
\usepackage{multirow}
\usepackage{wrapfig}

\title{DELTA: Edge-Centric Dual Attention for\\Compositional Reasoning on Knowledge Graphs}

\author{
  Benjamin Brown \\
  Independent Researcher \\
  \texttt{bdbrown4@eagles.usi.edu} \\
}

\begin{document}

\maketitle

\begin{abstract}
Graph neural networks treat edges as passive conduits for message passing between nodes, limiting their ability to reason about \emph{relationships between relationships}---the compositional structure underlying multi-hop inference in knowledge graphs. We introduce \textbf{DELTA} (\textbf{D}ual \textbf{E}dge-\textbf{L}inked \textbf{T}ransformer \textbf{A}rchitecture), which elevates edges to first-class computational entities via \emph{edge-to-edge attention}: edges attend to structurally adjacent edges through a multi-hop edge adjacency matrix, enabling direct relational composition without node-mediated bottlenecks. On three synthetic structural benchmarks (edge-type relation classification, noise robustness at 80\% corruption, and compositional relation discovery), evaluated across 5 seeds on an H100, DELTA outperforms GraphGPS and GRIT by $+0.573$, $+0.290$, and $+0.095$ accuracy respectively---providing the primary cross-task evidence that edge-to-edge attention captures a structurally distinct reasoning capacity. On the complete 14,541-entity FB15k-237 graph (mean degree $\approx$4.1), a 3-seed hop-depth ablation confirms that 2-hop edge adjacency improves 3-hop chain query MRR by $+0.012$ over 1-hop adjacency ($0.218$ vs.\ $0.206$), validating the core mechanism at realistic scale. As an exploratory characterization on a top-500-entity degree-based subgraph (mean degree $\approx$19.7), DELTA-Matched (158K parameters) is the only model among 7 tested with a depth-monotonic MRR improvement from 2-hop (0.758) to 5-hop (0.790); GraphGPS degrades from 0.754 to 0.690, and the advantage reaches $+0.100$ MRR at 5-hop. Cross-seed validation confirms the subgraph trend ($0.742 \pm 0.009$ vs.\ $0.713 \pm 0.007$ at 3-hop; non-overlapping CIs). However, a direct mechanism ablation on this dense subgraph finds 1-hop, 2-hop, and node-only conditions statistically indistinguishable, consistent with high link density saturating 1-hop adjacency coverage. An emergent edge--node temperature asymmetry arises during training: edge attention heads converge to sharper distributions while node heads prefer softer averaging, a pattern discovered autonomously via gradient descent. DELTA achieves this with 31\% fewer parameters than GraphGPS and faster per-query inference (0.8--0.9$\times$), at the cost of a one-time encoding overhead from edge adjacency construction. The central open question for future work is whether a SOTA-performing DELTA variant---one incorporating specialized KGC training strategies at full-graph scale---preserves the depth-monotonic compositional advantage observed here.
\end{abstract}

%----------------------------------------------------------------------
\section{Introduction}
\label{sec:intro}
%----------------------------------------------------------------------

Knowledge graph completion (KGC) requires inferring missing links from observed relational structure. While standard link prediction evaluates single-hop tail prediction (``1p''), real-world reasoning often demands \emph{compositional} multi-hop inference: chaining relations such as \texttt{born\_in} $\circ$ \texttt{located\_in} to derive \texttt{nationality}. This compositional capacity is precisely where existing architectures falter.

Embedding methods (TransE~\cite{bordes2013translating}, DistMult~\cite{yang2015embedding}, RotatE~\cite{sun2019rotate}) represent entities and relations as vectors without structural inductive bias. They excel at memorizing local link statistics but collapse on multi-hop chains: DistMult's MRR drops from 0.728 (2-hop) to 0.457 (5-hop) on our top-500-entity subgraph benchmark. Message-passing GNNs (CompGCN~\cite{vashishth2019composition}, R-GCN~\cite{schlichtkrull2018modeling}) and graph transformers (GraphGPS~\cite{rampavsek2022recipe}) propagate information through nodes, treating edges as scalar weights or simple feature carriers. This design creates an information bottleneck: to reason about relationships between relationships, information must route through intermediate nodes, accumulating compression loss at each hop.

We propose \textbf{DELTA}, an architecture that makes edges first-class computational entities. DELTA's key mechanism is \emph{edge-to-edge attention via multi-hop edge adjacency}: two edges that share a node are ``adjacent'' in edge space, and at 2 hops, an edge can attend to edges two steps away. This creates a direct channel for relational composition---\texttt{born\_in} can directly attend to \texttt{located\_in} without routing through the intermediate city node---enabling transitive inference as a natural consequence of the attention structure rather than an emergent property that must be learned.

DELTA runs node attention and edge attention as parallel streams, reconciled through a bidirectional bridge that ensures cross-stream coherence. This \emph{dual parallel attention} converges 2.7$\times$ faster than sequential alternatives while maintaining the same representational capacity.

\textbf{Contributions:}
\begin{enumerate}
    \item \textbf{Edge-to-edge attention architecture.} We introduce dual parallel attention with multi-hop edge adjacency, the first GNN mechanism that enables direct edge-to-edge reasoning without node-mediated bottlenecks (\S\ref{sec:architecture}).
    \item \textbf{Cross-task synthetic validation.} On three synthetic structural benchmarks---edge-type relation classification, edge classification under 80\% noise corruption, and compositional relation discovery (4 base $+$ 3 derived rules)---DELTA achieves 0.880, 1.000, and 1.000 mean accuracy (5 seeds, H100), vs.\ GraphGPS 0.293/0.697/0.905 and GRIT 0.307/0.710/0.893. These tasks are entirely distinct from knowledge graph completion and provide cross-task evidence that edge-to-edge attention captures a general structural reasoning capacity (\S\ref{sec:exp_synthetic}).
    \item \textbf{Full-scale validation on complete FB15k-237.} A direct 3-seed hop-depth ablation on the full 14,541-entity graph confirms that 2-hop edge adjacency improves 3-hop chain query MRR by $+0.012$ over 1-hop ($0.218$ vs.\ $0.206$) on the sparse graph (mean degree $\approx$4.1) where the mechanism is most relevant (\S\ref{sec:exp_ablation}).
    \item \textbf{Exploratory depth characterization on degree-based subgraph.} On a 500-entity degree-based subgraph of FB15k-237, DELTA-Matched exhibits a depth-monotonic MRR improvement from 2-hop to 5-hop that no other tested model shares, with the advantage over GraphGPS accelerating from $+0.004$ to $+0.100$ MRR. A direct mechanism ablation on this subgraph finds all hop-depth conditions statistically indistinguishable (consistent with high density saturating 1-hop adjacency), so these results characterize emergent behavioral depth-dependence rather than isolating the mechanism (\S\ref{sec:exp_compositional}).
    \item \textbf{Emergent attention temperature asymmetry.} We observe that edge and node attention heads develop fundamentally different sharpness preferences during training: edge heads drift toward higher temperature (sharper) while node heads drift lower (softer). This asymmetry is consistent across 20+ configurations and 9 experimental phases, suggesting distinct computational roles for the two attention streams (Appendix~\ref{app:temperature}).
    \item \textbf{Scalability characterization.} We evaluate DELTA at scales from $N{=}500$ to $N{=}14{,}541$ entities, identify attention dilution as the primary scaling bottleneck, and validate top-$k$ sparse attention ($k{=}128$) as a lossless approximation that preserves full-softmax quality (Appendix~\ref{app:scaling}).
\end{enumerate}

%----------------------------------------------------------------------
\section{Related Work}
\label{sec:related}
%----------------------------------------------------------------------

\paragraph{Knowledge graph embeddings.}
TransE~\cite{bordes2013translating} models relations as translations in embedding space; DistMult~\cite{yang2015embedding} uses bilinear scoring; RotatE~\cite{sun2019rotate} models relations as rotations in complex space; ComplEx~\cite{trouillon2016complex} extends to complex-valued embeddings. These methods achieve strong single-hop link prediction but lack structural inductive bias for multi-hop reasoning. PairRE~\cite{chao2021pairre} and RotatE achieve MRR $\sim$0.34 on full FB15k-237, but rely entirely on embedding geometry rather than graph structure for relational composition.

\paragraph{Message-passing GNNs for KGs.}
R-GCN~\cite{schlichtkrull2018modeling} introduces relation-specific weight matrices for typed edges. CompGCN~\cite{vashishth2019composition} jointly embeds nodes and relations through composition operations (subtraction, multiplication, circular correlation). Both treat edges as typed message channels rather than computational entities---edge representations are derived from relation type embeddings, not learned through edge-to-edge interaction. NBFNet~\cite{zhu2021neural} and RED-GNN~\cite{zhang2022knowledge} address multi-hop reasoning via path-based message passing but still route all information through nodes. A*Net~\cite{zhu2022astar} scales path-based reasoning to large KGs via priority-based entity selection. These path-based methods are a key baseline comparison target for future work; our primary comparison here is against general graph transformer architectures (GraphGPS, GRIT).

\paragraph{Graph transformers.}
GraphGPS~\cite{rampavsek2022recipe} combines local MPNN message passing with global transformer attention, using positional/structural encodings to inject graph awareness into the attention mechanism. GRIT~\cite{ma2023graph} uses random-walk structural encodings with graph transformers. Both architectures provide expressive node-level representations but lack mechanisms for direct edge-to-edge interaction. Exphormer~\cite{shirzad2023exphormer} addresses scalability via sparse attention patterns but does not introduce edge-level attention.

\paragraph{Edge-aware architectures.}
Line graphs~\cite{chen2017supervised} and their GNN variants operate on the dual graph where edges become nodes, enabling edge-level message passing. However, line graph construction is $O(E \cdot d_{\max})$ and the resulting graph can be much larger than the original. DELTA's edge adjacency matrix provides similar connectivity at $O(E^{0.97})$ via sparse COO operations, with the crucial difference that DELTA maintains both node and edge representations simultaneously through dual attention. EHGNN~\cite{gong2023ehgnn} incorporates edge features into heterogeneous GNNs but does not implement edge-to-edge attention.

\paragraph{Multi-hop reasoning on KGs.}
BetaE~\cite{ren2020beta} and ConE~\cite{zhang2021cone} embed logical queries as geometric regions for answering complex queries. GQE~\cite{hamilton2018embedding} and Q2B~\cite{ren2020query2box} project queries into embedding spaces. CQD~\cite{arakelyan2021cqd} answers complex queries by continuous optimization over link predictor scores. GNN-QE~\cite{zhu2022gnnqe} combines GNN encoders with fuzzy logic for complex query answering. These methods are designed specifically for complex query answering rather than general KGC. Our multi-hop evaluation uses chain queries (2p--5p) with soft entity traversal---a complementary protocol that tests compositional capacity within standard KGC training, and is not directly comparable to BetaE/GQE query benchmark numbers.

%----------------------------------------------------------------------
\section{Architecture}
\label{sec:architecture}
%----------------------------------------------------------------------

DELTA processes a knowledge graph $G = (V, E, R)$ with entity set $V$, edge set $E \subseteq V \times R \times V$, and relation set $R$. Each entity $v_i$ has a learned embedding $\mathbf{h}_i \in \mathbb{R}^{d_n}$ and each edge $(v_i, r_k, v_j)$ has a learned embedding $\mathbf{e}_{ij} \in \mathbb{R}^{d_e}$.

\subsection{Multi-Hop Edge Adjacency}
\label{sec:edge_adj}

The foundation of DELTA is the \emph{edge adjacency matrix} $\mathbf{A}_E$, which encodes structural connectivity between edges. Two edges $e_a = (v_i, r, v_j)$ and $e_b = (v_k, r', v_l)$ are 1-hop adjacent if they share an endpoint: $\{v_i, v_j\} \cap \{v_k, v_l\} \neq \emptyset$. The 2-hop edge adjacency extends this via the incidence structure:

\begin{equation}
    \mathbf{A}_E^{(2)} = (\mathbf{B}^\top \mathbf{B})^2 - \text{diag}
\end{equation}

where $\mathbf{B} \in \{0,1\}^{|V| \times |E|}$ is the node-edge incidence matrix. This is computed efficiently using sparse COO tensor operations via \texttt{torch\_sparse.spspmm}, achieving $O(E^{0.97})$ scaling (validated empirically from 50 to 2,500 edges). At 2 hops, an edge can ``see'' edges two steps away, enabling direct reasoning about relational chains such as \texttt{born\_in} $\to$ \texttt{located\_in} without intermediate node representations.

We validate the role of 2-hop edge adjacency in \S\ref{sec:exp_synthetic}: on a synthetic transitive relation benchmark, 2-hop achieves 100\% accuracy on derived relations vs.\ 61.1\% for 1-hop. On FB15k-237 multi-hop queries, a direct ablation finds all three adjacency configurations perform comparably (Table~\ref{tab:ablation_hop}); we discuss this result and its implications in \S\ref{sec:exp_ablation}.

\subsection{Edge-to-Edge Attention}
\label{sec:edge_attn}

Given the edge adjacency $\mathbf{A}_E^{(2)}$ stored as a sparse index $[2, E_{\text{adj}}]$, we compute multi-head edge attention. For edge $e_i$ attending to adjacent edge $e_j$:

\begin{equation}
    \alpha_{ij}^{(h)} = \frac{\tau^{(h)} \cdot (\mathbf{W}_Q^{(h)} \mathbf{e}_i)^\top (\mathbf{W}_K^{(h)} \mathbf{e}_j)}{\sqrt{d_h}} + \mathbf{W}_{\text{ctx}}[\mathbf{h}_{s(j)} \| \mathbf{h}_{t(j)}]_h
    \label{eq:edge_attn}
\end{equation}

where $\tau^{(h)} = \exp(\theta^{(h)})$ is a learnable per-head temperature (always positive). Note that $\tau^{(h)}$ \emph{multiplies} the dot product (not divides): larger $\tau$ produces sharper, more selective attention---the inverse of the conventional divide-by-temperature softmax formulation. $s(j)$ and $t(j)$ are the source and target nodes of edge $e_j$, $\|$ denotes concatenation, and $d_h = d_e / H$ is the head dimension. The node context bias $\mathbf{W}_{\text{ctx}} \in \mathbb{R}^{2d_n \times H}$ grounds edge attention in the endpoint identities of the attending edges.

Attention weights are computed via scatter-softmax over the sparse adjacency structure (normalizing per target edge), followed by weighted aggregation:

\begin{equation}
    \mathbf{e}_i' = \text{LayerNorm}\left(\mathbf{e}_i + \mathbf{W}_O \sum_{j \in \mathcal{N}_E(i)} \text{softmax}(\alpha_{ij}) \cdot \mathbf{W}_V \mathbf{e}_j\right)
\end{equation}

\textbf{Top-$k$ sparse attention.} For large graphs where $|\mathcal{N}_E(i)|$ becomes prohibitive, we retain only the $k$ highest-scoring neighbors per target edge (based on mean attention across heads), then renormalize. At $N{=}5{,}000$ entities with 63M adjacency pairs, $k{=}128$ preserves full-softmax quality exactly (test MRR $0.2472$ vs.\ $0.2457$ for full softmax), while $k{=}64$ degrades by $-5.6\%$.

\subsection{Node Attention}
\label{sec:node_attn}

Node attention follows a GAT-style~\cite{velickovic2018graph} design where each node attends to its neighbors, with edge features providing a per-head bias:

\begin{equation}
    \beta_{ij}^{(h)} = \frac{\tau_n^{(h)} \cdot (\mathbf{W}_Q^{(h)} \mathbf{h}_i)^\top (\mathbf{W}_K^{(h)} \mathbf{h}_j)}{\sqrt{d_h}} + \mathbf{W}_{\text{edge}}(\mathbf{e}_{ij})_h
\end{equation}

Node attention uses independent learnable temperatures $\tau_n^{(h)}$ per head, enabling the model to discover optimal sharpness for each attention stream independently.

\subsection{Dual Parallel Attention}
\label{sec:dual_attn}

Both attention mechanisms operate on the \emph{same} input state simultaneously:

\begin{equation}
    \mathbf{H}' = \text{NodeAttn}(\mathbf{H}, \mathbf{E}, G), \qquad \mathbf{E}' = \text{EdgeAttn}(\mathbf{E}, \mathbf{H}, \mathbf{A}_E)
\end{equation}

This parallel design ensures that neither stream's output biases the other's computation within a layer, and converges 2.7$\times$ faster than sequential alternatives.

\subsection{Reconciliation Bridge}
\label{sec:reconciliation}

After parallel attention, a bidirectional bridge ensures cross-stream coherence:

\begin{align}
    \hat{\mathbf{e}}_{ij} &= \text{LN}(\mathbf{e}'_{ij} + \mathbf{W}_1[\mathbf{e}'_{ij} \| \mathbf{h}'_{s(ij)} \| \mathbf{h}'_{t(ij)}]) & \text{(edges absorb node context)} \\
    \hat{\mathbf{h}}_i &= \text{LN}(\mathbf{h}'_i + \mathbf{W}_2[\mathbf{h}'_i \| \bar{\mathbf{e}}_i]) & \text{(nodes absorb edge context)}
\end{align}

where $\bar{\mathbf{e}}_i = \text{mean}(\{\hat{\mathbf{e}}_{ij} : j \in \mathcal{N}(i)\})$ aggregates incident updated edge features. Both updates use residual connections with LayerNorm. Edges update first (absorbing node context), then nodes update (absorbing the already node-informed edge features), creating a sequential cascade within the bridge.

\subsection{Link Prediction Head}

For link prediction, we use a DistMult~\cite{yang2015embedding} decoder:

\begin{equation}
    s(v_i, r_k, v_j) = \mathbf{h}_i^\top \text{diag}(\mathbf{r}_k) \mathbf{h}_j
\end{equation}

where $\mathbf{r}_k$ is a learned relation embedding. Training uses 1-vs-all binary cross-entropy with label smoothing (0.1), following standard KGC practice.

\subsection{Multi-Hop Scoring}

For $n$-hop path queries $(v_0, r_1, ?, r_2, ?, \ldots, r_n, ?)$, we use \emph{soft entity traversal}: at each intermediate hop, score all entities via the DistMult decoder, apply softmax with temperature $\tau=1.0$ to obtain a probability distribution over entities, take the weighted average entity embedding, and pass it to the next hop. This provides a differentiable approximation to hard entity traversal while remaining faithful to the compositional structure of the query.

%----------------------------------------------------------------------
\section{Experiments}
\label{sec:experiments}
%----------------------------------------------------------------------

\subsection{Setup}
\label{sec:exp_setup}

\paragraph{Dataset.} We use FB15k-237~\cite{toutanova2015representing}, the standard filtered subset of Freebase with 14,541 entities, 237 relations, and 310,116 triples. Our compositional reasoning and mechanism analysis experiments (\S\ref{sec:exp_compositional}) use a \textbf{top-500 degree subgraph} (494 entities, 207 relations, ${\sim}9{,}700$ train triples, mean degree $\approx$19.7) with full filtered evaluation. This degree-biased subgraph enables rapid multi-configuration ablations across many seeds and depth levels. Scaling experiments (Appendix~\ref{app:scaling}) extend to $N{=}2{,}000$ and $N{=}5{,}000$ entity subgraphs. The hop-depth ablation (\S\ref{sec:exp_ablation}) additionally reports a full-scale evaluation on the complete 14,541-entity graph (mean degree $\approx$4.1, 310,116 triples) to validate key findings at standard benchmark scale.

\paragraph{Models.} We compare seven configurations:
\begin{itemize}
    \item \textbf{DELTA-Matched} (158K params): 1 DELTA layer, $d_n{=}64$, $d_e{=}32$, 4 heads. Our primary model, designed for parameter-fair comparison.
    \item \textbf{DELTA-Full} (293K params): 3 DELTA layers.
    \item \textbf{SelfBootstrapHybrid} (381K params): A FixedChain DELTA layer bootstraps graph construction before 3 DELTA processing layers.
    \item \textbf{GraphGPS}~\cite{rampavsek2022recipe} (228K params): Local MPNN + global transformer attention.
    \item \textbf{GRIT}~\cite{ma2023graph} (197K params): Random-walk structural encodings + graph transformer.
    \item \textbf{DistMult}~\cite{yang2015embedding} (47K params): Pure embedding baseline (no GNN).
    \item \textbf{SelfBootstrap} (299K params): DELTA bootstraps its own graph construction.
\end{itemize}

\paragraph{Training.} All models use Adam optimizer ($\text{lr}=0.003$), batch size 4,096, 1-vs-all BCE loss with label smoothing 0.1, and early stopping with patience 10. Evaluation uses standard filtered MRR and Hits@$k$ metrics. Multi-hop queries use 10\% DropEdge for DELTA (validated in \S\ref{sec:exp_compositional}).

\paragraph{Multi-hop query generation.} Chain queries are generated by composing training edges to reach test edges: a 2-hop query chains one train edge to one test edge, a 3-hop query chains two train edges to one test edge, etc. All queries undergo strict leakage prevention: no 1-hop shortcuts between anchor and answer in the training graph, no duplicate queries, and automated audit verification before evaluation. Query counts: 1p=486, 2p=5,382, 3p=10,000, 4p=10,000, 5p=10,000.

%----------------------------------------------------------------------
\subsection{Synthetic Structural Benchmarks}
\label{sec:exp_synthetic}
%----------------------------------------------------------------------

Before evaluating on the FB15k-237 knowledge graph, we validate the edge-to-edge attention mechanism on synthetic structural tasks that isolate compositional reasoning without confounds from embedding quality or training scale.

\paragraph{Transitive edge adjacency (Phase 11).}
On a 100-entity synthetic knowledge graph with explicitly transitive relations, we ablate edge adjacency depth (Table~\ref{tab:phase11}). 2-hop edge adjacency achieves \textbf{100\%} accuracy on derived (transitive) relation instances, vs.\ 1-hop adjacency at 61.1\%, a node GNN at 83.3\%, and an MLP at 37.5\%. 2-hop adjacency is necessary and sufficient for transitive relational composition. This result, alongside the full-graph ablation ($+0.012$ on 3p MRR), constitutes the primary mechanistic evidence for the hop-depth claim, since the dense-subgraph mechanism ablation (Phase 66, Table~\ref{tab:ablation_hop}) finds all conditions statistically indistinguishable.

\begin{table}[t]
\centering
\caption{Phase 11 --- transitive relation classification accuracy by adjacency configuration (100-entity synthetic graph, single seed). Only 2-hop edge adjacency achieves perfect accuracy on derived (composed) relation instances.}
\label{tab:phase11}
\begin{tabular}{@{}lcc@{}}
\toprule
\textbf{Config} & \textbf{Transitive Acc.} & \textbf{$\Delta$ vs.\ 2-hop} \\
\midrule
2-hop (DELTA) & \textbf{1.000} & --- \\
Node-GNN (no edge attn) & 0.833 & $-0.167$ \\
1-hop & 0.611 & $-0.389$ \\
MLP (no structure) & 0.375 & $-0.625$ \\
\bottomrule
\end{tabular}
\end{table}

\paragraph{Three-task benchmark (Phase 34).}
We compare DELTA-Matched, GraphGPS, and GRIT on three synthetic structural tasks (5 seeds $\times$ 500 epochs, H100; Table~\ref{tab:synthetic}):
\begin{itemize}
    \item \textbf{Edge-type classification:} A synthetic KG with 6 relation patterns (8 instances each); classify each edge by its relation type. DELTA achieves 0.880 vs.\ GraphGPS 0.293 and GRIT 0.307 (advantage $+0.573$). Edge-to-edge attention learns discriminative edge-type representations that node-centric message passing cannot extract---standard graph transformers have no native edge-classification head, and their node representations lose the information needed to distinguish edge types.
    \item \textbf{Noise robustness:} Repeat edge-type classification with 80\% of edges randomly corrupted. DELTA maintains 1.000 accuracy while GraphGPS degrades to 0.697 and GRIT to 0.710. DELTA's edge-level representation is structurally grounded and robust to corrupted neighborhood signals.
    \item \textbf{Compositional relation discovery:} A KG with 4 base relation types and 3 derived (compositional) rules; models must classify edges including derived instances produced by relational composition. DELTA achieves 1.000 vs.\ GraphGPS 0.905 ($+0.095$), directly demonstrating that 2-hop edge adjacency enables compositional rule discovery without supervision.
\end{itemize}

\begin{table}[t]
\centering
\caption{Synthetic structural benchmarks (Phase 34): mean accuracy across 5 seeds (500 epochs, H100), compared at matched parameter budgets. Advantage row shows DELTA vs.\ the stronger of GraphGPS/GRIT.}
\label{tab:synthetic}
\begin{tabular}{@{}lcccc@{}}
\toprule
\textbf{Model} & \textbf{Params} & \textbf{Edge Cls.} & \textbf{Noise @0.8} & \textbf{Path Comp.} \\
\midrule
\textbf{DELTA-Matched} & 158K & \textbf{0.880} & \textbf{1.000} & \textbf{1.000} \\
GraphGPS & 228K & 0.293 & 0.697 & 0.905 \\
GRIT & 197K & 0.307 & 0.710 & 0.893 \\
\midrule
\textit{DELTA advantage} & --- & \textbf{+0.573} & \textbf{+0.290} & \textbf{+0.095} \\
\bottomrule
\end{tabular}
\end{table}

The cross-task pattern is consistent: DELTA's largest advantage is on the edge-type classification task ($+0.573$), where node-centric architectures are structurally mismatched. The advantage is smallest on path composition ($+0.095$), where multi-hop node-level message passing provides a partial substitute. Together with the transitive adjacency ablation, these results show that DELTA's mechanism is not FB15k-237-specific---edge-to-edge attention provides a general structural advantage wherever edge identity matters.

\textbf{Probe architecture note.} GraphGPS/GRIT have no native edge-classification head; their probes use mean-pooled endpoint node embeddings, while DELTA's probe operates on learned edge representations. This asymmetry may inflate DELTA's advantage on edge-type classification. The path composition result ($+0.095$), where all models share the same endpoint prediction protocol, provides a cleaner cross-architecture comparison.

%----------------------------------------------------------------------
\subsection{Compositional Reasoning}
\label{sec:exp_compositional}
%----------------------------------------------------------------------

We characterize depth-dependent performance on the FB15k-237 degree-based subgraph. DELTA-Matched is the only model among 7 tested that shows a depth-monotonic MRR improvement from 2-hop to 5-hop; GraphGPS and all other baselines degrade. Readers should note that a direct mechanism ablation on the same subgraph (Table~\ref{tab:ablation_hop}) finds 1-hop, 2-hop, and node-only conditions statistically indistinguishable---the depth-monotonic pattern is a behavioral property of the model on this graph, not an isolated consequence of the hop-depth setting. The full-graph ablation (\S\ref{sec:exp_ablation}) provides the mechanistic evidence that 2-hop adjacency adds discriminative value on sparse graphs.

\begin{table}[t]
\centering
\caption{Multi-hop MRR by reasoning depth (single seed, 0\% DropEdge). DELTA-Matched is the only model with higher MRR at 5-hop than 2-hop (net $+0.032$). All baselines show net degradation from 2-hop to 5-hop.}
\label{tab:depth}
\begin{tabular}{@{}lcccccc@{}}
\toprule
\textbf{Model} & \textbf{Params} & \textbf{1p} & \textbf{2p} & \textbf{3p} & \textbf{4p} & \textbf{5p} \\
\midrule
\textbf{DELTA-Matched} & 158K & 0.541 & \textbf{0.758} & \textbf{0.753} & \textbf{0.767} & \textbf{0.790} \\
GraphGPS & 228K & \textbf{0.523} & 0.754 & 0.727 & 0.701 & 0.690 \\
DistMult & 47K & 0.494 & 0.728 & 0.583 & 0.511 & 0.457 \\
\midrule
\multicolumn{2}{@{}l}{\textit{DELTA -- GraphGPS gap}} & +0.018 & +0.004 & +0.026 & +0.066 & \textbf{+0.100} \\
\bottomrule
\end{tabular}
\end{table}

\paragraph{The depth advantage accelerates.}
Table~\ref{tab:depth} presents the central result. DELTA-Matched's lead over GraphGPS grows from $+0.004$ (2p) to $+0.100$ (5p)---a 25$\times$ amplification. Critically, DELTA's 5p MRR (0.790) \emph{exceeds} its own 2p MRR (0.758). No other model exhibits this property: GraphGPS loses $-0.064$ MRR from 2p$\to$5p, while DistMult loses $-0.271$.

\paragraph{Consistent multi-seed advantage.}
Table~\ref{tab:multiseed} confirms statistical robustness (3 seeds, 10\% DropEdge): DELTA improves from 2p ($0.730 \pm 0.011$) to 3p ($0.742 \pm 0.009$), while GraphGPS degrades ($0.727\to 0.713$); confidence intervals do not overlap. The full 7-model single-seed comparison (Appendix~\ref{app:multihop_full}, Table~\ref{tab:all_models}) shows DELTA-Matched is the only model with a positive 2p$\to$3p delta ($+0.005$); every other model degrades.


\begin{table}[t]
\centering
\caption{Multi-seed 3-hop MRR ($\pm$ std, 3 seeds). Non-overlapping confidence intervals.}
\label{tab:multiseed}
\begin{tabular}{@{}lccccc@{}}
\toprule
\textbf{Config} & \textbf{Params} & \textbf{1p} & \textbf{2p} & \textbf{3p} & \textbf{2p$\to$3p} \\
\midrule
DELTA @10\% DE & 158K & $0.543 \pm 0.006$ & $0.730 \pm 0.011$ & $\mathbf{0.742 \pm 0.009}$ & +0.012 \\
GraphGPS & 228K & $0.529 \pm 0.009$ & $0.727 \pm 0.007$ & $0.713 \pm 0.007$ & $-0.014$ \\
\bottomrule
\end{tabular}
\end{table}

\paragraph{Robustness to regularization.}
We sweep DropEdge rates (0\%--40\%) for both DELTA-Matched and GraphGPS (Appendix~\ref{app:dropedge}). DELTA leads GraphGPS on 3p MRR at all five rates tested, with the advantage ranging from $+0.001$ (at 20\%, within noise) to $+0.029$ (at 0\%). The advantage is statistically meaningful at four of the five rates; the multi-hop advantage is structural, not an artifact of a particular hyperparameter configuration.

%----------------------------------------------------------------------
\subsection{Standard Link Prediction}
\label{sec:exp_lp}
%----------------------------------------------------------------------

On standard 1p link prediction (top-500 subgraph), GraphGPS leads DELTA-Matched by $0.018$ MRR (0.513 vs.\ 0.495). Full per-model results in Appendix~\ref{app:lp_full} (Table~\ref{tab:lp}).

%----------------------------------------------------------------------
\subsection{Ablation: 1-Hop vs.\ 2-Hop Edge Adjacency}
\label{sec:exp_ablation}
%----------------------------------------------------------------------

To directly validate the role of 2-hop edge adjacency---the architectural mechanism underlying DELTA's compositional reasoning---we ablate three configurations on the FB15k-237 subgraph using the multi-hop evaluation protocol:

\begin{itemize}
    \item \textbf{1-hop}: edges attend only to directly adjacent edges (share a node)
    \item \textbf{2-hop}: edges attend to edges two steps away (default DELTA)
    \item \textbf{Node-only}: edge attention stream disabled; DELTA reduces to GAT-style node attention
\end{itemize}

\begin{table}[t]
\centering
\caption{Hop-depth ablation (Phase 66): DELTA-Matched on FB15k-237 top-500 subgraph, 3 seeds $\times$ 500 epochs. Values are mean $\pm$ std. All differences are within 1$\sigma$; no condition is statistically superior.}
\label{tab:ablation_hop}
\begin{tabular}{@{}lccccc@{}}
\toprule
\textbf{Config} & \textbf{LP MRR} & \textbf{1p MRR} & \textbf{2p MRR} & \textbf{3p MRR} & \textbf{2p$\to$3p} \\
\midrule
2-hop (DELTA)  & $0.496\pm0.003$ & $0.534\pm0.008$ & $0.726\pm0.014$ & $0.731\pm0.006$ & $+0.005$ \\
1-hop          & $0.498\pm0.008$ & $0.551\pm0.001$ & $0.721\pm0.007$ & $0.729\pm0.007$ & $+0.008$ \\
Node-only      & $0.505\pm0.006$ & $0.553\pm0.006$ & $0.728\pm0.011$ & $0.742\pm0.012$ & $+0.014$ \\
\bottomrule
\end{tabular}
\end{table}

Table~\ref{tab:ablation_hop} shows all three configurations are statistically indistinguishable on the dense subgraph (every pairwise gap within one standard deviation). The node-only condition achieves the highest point estimates, though within noise---consistent with 1-hop adjacency already saturating the structural neighborhood at mean degree $\approx 19.7$. This motivates the full-graph ablation below, where 1-hop neighborhoods are genuinely sparse.

\paragraph{Full-Scale Ablation (Phases 67--68).}
On the complete 14,541-entity graph (mean degree $\approx 4.1$, top-$k{=}128$ sparse attention, 1,064K parameters, 3 seeds), hops=2 achieves $\mathbf{+0.012}$ on 3p MRR over hops=1, exceeding the pre-specified $+0.010$ threshold. (Preliminary 2-seed results showed $+0.017$; seed~456's hops=2 3p score revised the 3-seed average to $+0.012$, still above threshold.) LP MRR is statistically equivalent ($0.204$ vs.\ $0.205$). On a random 2,242-entity subgraph (uniformly sampled, mean degree~7.6, equal adjacency budget), the 3p gap is $+0.010$ over 3 seeds. Full multi-metric breakdowns in Appendix~\ref{app:ablation_full} (Table~\ref{tab:phase67} Phase~67 and Table~\ref{tab:phase68} Phase~68). Table~\ref{tab:cross_density} synthesizes all three density regimes.

\begin{table}[h]
\small
\centering
\caption{Cross-density summary: hops=2 vs.\ hops=1 on 3p MRR (Phases 66--68). The gap decreases monotonically with mean degree. Absolute 3p values are not directly comparable across rows.}
\label{tab:cross_density}
\resizebox{\linewidth}{!}{\begin{tabular}{@{}llcccc@{}}
\toprule
\textbf{Phase} & \textbf{Subgraph} & \textbf{Mean Deg} & \textbf{hops=1\ 3p} & \textbf{hops=2\ 3p} & \textbf{$\Delta$\ 3p} \\
\midrule
67 & Full FB15k-237 ($N{=}14{,}541$) & 4.1  & $0.206\pm0.005$ & $0.218\pm0.008$ & $\mathbf{+0.012}$ \\
68 & Random ($N{=}2{,}242$)          & 7.6  & $0.274\pm0.060$ & $0.284\pm0.031$ & $+0.010$ \\
66 & Top-degree ($N{=}500$)          & 19.7 & $0.729\pm0.007$ & $0.731\pm0.006$ & $+0.002$ \\
\bottomrule
\end{tabular}}
\end{table}

Across all three density regimes the hops=2 advantage decreases monotonically with mean degree, consistent with 1-hop adjacency becoming informationally sufficient above a density threshold.

For inference timing, scaling analysis, and attention temperature mechanistic details, see Appendix~\ref{app:inference}--\ref{app:temperature}.

%----------------------------------------------------------------------
\section{Discussion}
\label{sec:discussion}
%----------------------------------------------------------------------

\paragraph{Why does the advantage grow with depth?}
We hypothesize that 2-hop edge adjacency creates a multiplicative compositional advantage: DELTA composes relational information directly in edge space, while node-based architectures decompress through node representations at each hop. At 1 hop this overhead is negligible; at 5 hops, accumulated compression loss in node-mediated architectures becomes substantial. The temperature asymmetry (Appendix~\ref{app:temperature})---edge attention sharper than node attention---supports this interpretation: the model autonomously discovers differential sharpness that mirrors the two streams' distinct computational roles.

\paragraph{When does the mechanism matter?}
The Phase 66 dense-subgraph ablation (mean degree ${\approx}19.7$) finds all adjacency conditions statistically indistinguishable, while the Phase 67 full-graph ablation (mean degree ${\approx}4.1$) finds $+0.012$ on 3p MRR. The mechanism matters when graph sparsity creates relational chains that 1-hop adjacency cannot bridge. Table~\ref{tab:cross_density} confirms this monotonically across all three density regimes.

\paragraph{Does the compositional inductive bias trade off against 1p performance?}
DELTA-Matched (158K parameters) consistently outperforms DELTA-Full (293K) on multi-hop tasks, which we interpret as larger models memorizing local link statistics at the expense of compositional generalization. The full-graph LP gap below SOTA ($0.205$ vs.\ $0.338$--$0.415$) likely reflects training-strategy differences---self-adversarial negative sampling, longer schedules, entity-specific regularization---rather than an architectural incompatibility. The per-scale gap on the top-500 subgraph is only $-0.018$ against GraphGPS (97\% of its MRR at 69\% of parameters), suggesting the mechanism itself does not preclude competitive 1p performance; optimizing for it is an engineering question, not a structural limitation.

%----------------------------------------------------------------------
\section{Conclusion}
\label{sec:conclusion}
%----------------------------------------------------------------------

We have presented DELTA, an edge-centric architecture using edge-to-edge attention over multi-hop adjacency. A consistent picture emerges across three experimental contexts: edge-to-edge attention provides a structural reasoning advantage that grows with compositional depth and matters most when graph sparsity creates relational chains that 1-hop adjacency cannot bridge. On synthetic benchmarks, DELTA outperforms GraphGPS and GRIT by $+0.573$, $+0.290$, $+0.095$ (5 seeds, cross-task). On complete FB15k-237 (3 seeds), 2-hop adjacency yields $+0.012$ on 3p MRR over 1-hop, confirming the mechanism at realistic scale. On the degree-based subgraph (7 models, 5 depths), DELTA-Matched is the only model with depth-monotonic MRR improvement (GraphGPS: $-0.064$), while a direct ablation on the same dense subgraph finds conditions statistically indistinguishable---consistent with the mechanism mattering most when 1-hop adjacency cannot saturate the structural neighborhood. An emergent edge--node temperature asymmetry provides mechanistic interpretability.

These results suggest that the standard single-hop link prediction benchmark undervalues architectures with strong relational inductive biases. Multi-hop evaluation provides a more discriminating test of compositional capacity, and we encourage its adoption alongside standard metrics. We note clearly that DELTA's full-graph LP MRR ($\approx 0.205$) falls below published SOTA methods ($\approx 0.34$--$0.41$), reflecting the gap between demonstrating a novel mechanism and achieving optimized end-to-end performance---a gap that remains an important open problem.

\textbf{Future work} includes (1) full-depth evaluation (2-hop$\to$5-hop) on complete FB15k-237, comparison against NBFNet~\cite{zhu2021neural} and A*Net~\cite{zhu2022astar}, and BetaE/GQE benchmarks~\cite{ren2020beta}; and (2) cross-dataset validation (YAGO3-10, WN18RR, ogbl-wikikg2) and closing the SOTA LP gap via specialized KGC training while preserving the compositional advantage.

%----------------------------------------------------------------------
% References
%----------------------------------------------------------------------

\bibliographystyle{plain}

\begin{thebibliography}{30}

\bibitem{bordes2013translating}
Antoine Bordes, Nicolas Usunier, Alberto Garcia-Duran, Jason Weston, and Oksana Yakhnenko.
\newblock Translating embeddings for modeling multi-relational data.
\newblock In \emph{NeurIPS}, 2013.

\bibitem{yang2015embedding}
Bishan Yang, Wen-tau Yih, Xiaodong He, Jianfeng Gao, and Li Deng.
\newblock Embedding entities and relations for learning and inference in knowledge bases.
\newblock In \emph{ICLR}, 2015.

\bibitem{sun2019rotate}
Zhiqing Sun, Zhi-Hong Deng, Jian-Yun Nie, and Jian Tang.
\newblock RotatE: Knowledge graph embedding by relational rotation in complex space.
\newblock In \emph{ICLR}, 2019.

\bibitem{trouillon2016complex}
Th{\'e}o Trouillon, Johannes Welbl, Sebastian Riedel, {\'E}ric Gaussier, and Guillaume Bouchard.
\newblock Complex embeddings for simple link prediction.
\newblock In \emph{ICML}, 2016.

\bibitem{chao2021pairre}
Linlin Chao, Jianshan He, Taifeng Wang, and Wei Chu.
\newblock PairRE: Knowledge graph embeddings via paired relation vectors.
\newblock In \emph{ACL}, 2021.

\bibitem{schlichtkrull2018modeling}
Michael Schlichtkrull, Thomas~N. Kipf, Peter Bloem, Rianne van~den Berg, Ivan Titov, and Max Welling.
\newblock Modeling relational data with graph convolutional networks.
\newblock In \emph{ESWC}, 2018.

\bibitem{vashishth2019composition}
Shikhar Vashishth, Soumya Sanyal, Vikram Nitin, and Partha Talukdar.
\newblock Composition-based multi-relational graph convolutional networks.
\newblock In \emph{ICLR}, 2020.

\bibitem{zhu2021neural}
Zhaocheng Zhu, Zuobai Zhang, Louis-Pascal Xhonneux, and Jian Tang.
\newblock Neural Bellman-Ford networks: A general graph neural network framework for link prediction.
\newblock In \emph{NeurIPS}, 2021.

\bibitem{zhang2022knowledge}
Yongqi Zhang and Quanming Yao.
\newblock Knowledge graph reasoning with relational digraph.
\newblock In \emph{NeurIPS}, 2022.

\bibitem{rampavsek2022recipe}
Ladislav Ramp{\'a}{\v{s}}ek, Michael Galkin, Vijay Prakash Dwivedi, Anh~Tuan Luu, Guy Wolf, and Dominique Beaini.
\newblock Recipe for a general, powerful, scalable graph transformer.
\newblock In \emph{NeurIPS}, 2022.

\bibitem{ma2023graph}
Liheng Ma, Chen Lin, Derek Lim, Adriana Romero-Soriano, Puneet~K. Dokania, Mark Coates, Philip Torr, and Ser-Nam Lim.
\newblock Graph inductive biases in transformers without message passing.
\newblock In \emph{ICML}, 2023.

\bibitem{shirzad2023exphormer}
Hamed Shirzad, Ameya Velingker, Balaji Venkatachalam, Danica~J. Sutherland, and Ali~Kemal Sinop.
\newblock Exphormer: Sparse transformers for graphs.
\newblock In \emph{ICML}, 2023.

\bibitem{chen2017supervised}
Zhengdao Chen, Xiang Li, and Joan Bruna.
\newblock Supervised community detection with line graph neural networks.
\newblock In \emph{ICLR}, 2019.

\bibitem{gong2023ehgnn}
Jiaqi Gong, Shichao Pei, Jiaxin Bai, and Xiangliang Zhang.
\newblock EHGNN: Edge heterogeneous graph neural network.
\newblock \emph{arXiv preprint arXiv:2310.07242}, 2023.

\bibitem{ren2020beta}
Hongyu Ren and Jure Leskovec.
\newblock Beta embeddings for multi-hop logical reasoning in knowledge graphs.
\newblock In \emph{NeurIPS}, 2020.

\bibitem{zhang2021cone}
Zhanqiu Zhang, Jie Wang, Jiajun Chen, Shuiwang Ji, and Feng Wu.
\newblock ConE: Cone embeddings for multi-hop reasoning over knowledge graphs.
\newblock In \emph{NeurIPS}, 2021.

\bibitem{hamilton2018embedding}
William~L. Hamilton, Payal Bajaj, Marinka Zitnik, Dan Jurafsky, and Jure Leskovec.
\newblock Embedding logical queries on knowledge graphs.
\newblock In \emph{NeurIPS}, 2018.

\bibitem{ren2020query2box}
Hongyu Ren, Weihua Hu, and Jure Leskovec.
\newblock Query2Box: Reasoning over knowledge graphs in vector space using box embeddings.
\newblock In \emph{ICLR}, 2020.

\bibitem{toutanova2015representing}
Kristina Toutanova and Danqi Chen.
\newblock Observed versus latent features for knowledge base and text inference.
\newblock In \emph{3rd Workshop on Continuous Vector Space Models and their Compositionality}, 2015.

\bibitem{velickovic2018graph}
Petar Veli{\v{c}}kovi{\'c}, Guillem Cucurull, Arantxa Casanova, Adriana Romero, Pietro Li{\`o}, and Yoshua Bengio.
\newblock Graph attention networks.
\newblock In \emph{ICLR}, 2018.

\bibitem{zhu2022astar}
Zhaocheng Zhu, Xinyu Yuan, Michael Galkin, Sophie Xhonneux, Ming Zhang, Maxime Gazeau, and Jian Tang.
\newblock A*Net: A scalable path-based reasoning approach for knowledge graphs.
\newblock In \emph{NeurIPS}, 2023.


\bibitem{arakelyan2021cqd}
Erik Arakelyan, Daniel Daza, Pasquale Minervini, and Michael Cochez.
\newblock Complex query answering with neural link predictors.
\newblock In \emph{ICLR}, 2021.

\bibitem{zhu2022gnnqe}
Zhaocheng Zhu, Mikhail Galkin, Zuobai Zhang, and Jian Tang.
\newblock Neural-symbolic models for logical queries on knowledge graphs.
\newblock In \emph{ICML}, 2022.

\bibitem{dettmers2018convolutional}
Tim Dettmers, Pasquale Minervini, Pontus Stenetorp, and Sebastian Riedel.
\newblock Convolutional 2D knowledge graph embeddings.
\newblock In \emph{AAAI}, 2018.

\end{thebibliography}

%----------------------------------------------------------------------
% Appendix
%----------------------------------------------------------------------
\newpage
\appendix

\section{Hits@10 by Depth}
\label{app:hits10}

\begin{table}[h]
\centering
\caption{Hits@10 across reasoning depths (single seed, same experiment as Table~\ref{tab:depth}).}
\label{tab:hits10_depth}
\begin{tabular}{@{}lccccc@{}}
\toprule
\textbf{Model} & \textbf{1p} & \textbf{2p} & \textbf{3p} & \textbf{4p} & \textbf{5p} \\
\midrule
DELTA-Matched & 0.854 & 0.889 & \textbf{0.881} & \textbf{0.885} & \textbf{0.895} \\
GraphGPS & \textbf{0.848} & \textbf{0.897} & 0.859 & 0.818 & 0.826 \\
DistMult & 0.800 & 0.858 & 0.744 & 0.672 & 0.587 \\
\bottomrule
\end{tabular}
\end{table}

\section{DropEdge Robustness}
\label{app:dropedge}

\begin{table}[h]
\centering
\caption{3-hop MRR across DropEdge rates. DELTA-Matched leads at all rates.}
\label{tab:dropedge}
\begin{tabular}{@{}lccccc@{}}
\toprule
\textbf{Model} & \textbf{0\%} & \textbf{10\%} & \textbf{20\%} & \textbf{30\%} & \textbf{40\%} \\
\midrule
DELTA-Matched & 0.740 & 0.744 & 0.724 & 0.732 & 0.744 \\
GraphGPS & 0.711 & 0.716 & 0.723 & 0.720 & 0.725 \\
\midrule
$\Delta$ & +0.029 & +0.028 & +0.001 & +0.012 & +0.019 \\
\bottomrule
\end{tabular}
\end{table}

\section{Brain Architecture (Preliminary)}
\label{app:brain}

The BrainEncoder module uses Gumbel-sigmoid differentiable edge selection to construct new edges from learned node representations, then processes the augmented graph (original + constructed edges) through DELTA layers. On the FB15k-237 subgraph ($N{=}494$):

\begin{table}[h]
\centering
\caption{BrainEncoder results (multi-seed, density $d{=}0.01$, 200 epochs).}
\begin{tabular}{@{}lcccc@{}}
\toprule
\textbf{Config} & \textbf{MRR} & \textbf{H@1} & \textbf{H@10} & \textbf{Constructed edges} \\
\midrule
DELTA-Full & 0.494 & 0.355 & 0.760 & 0 \\
brain\_hybrid ($d$=0.01) & $0.484 \pm 0.010$ & $0.334 \pm 0.011$ & $\mathbf{0.799 \pm 0.006}$ & 2,435 \\
\bottomrule
\end{tabular}
\end{table}

The brain\_hybrid configuration matches DELTA-Full on MRR while delivering $+4.7\%$ Hits@10, demonstrating that differentiable graph construction adds genuine recall through learned structural augmentation. Density is a critical hyperparameter: $d{=}0.01$ (2,435 edges) strictly dominates $d{=}0.02$ (4,870 edges) on all metrics.

\section{Domain Transfer (Preliminary)}
\label{app:transfer}

A frozen DELTA encoder trained on FB15k-237, evaluated via a linear probe on WN18RR, achieves accuracy of \textbf{0.961}---without any domain adaptation mechanism (no GRL, no fine-tuning). \textbf{Note:} this evaluation uses 100 training samples, giving approximate $\pm$5\% confidence intervals; results should be treated as a preliminary signal, not a statistically validated finding. A proper transfer evaluation requires $\geq$10,000 samples with confidence intervals on held-out test splits. We report this result as a directional indicator that DELTA's learned edge representations may capture transferable relational structure.

%----------------------------------------------------------------------
\section{Full Multi-Hop and Link Prediction Results}
\label{app:multihop_full}
%----------------------------------------------------------------------

\begin{table}[h]
\centering
\caption{Complete multi-hop results (7 models, seed=1, 0\% DropEdge). Only DELTA-Matched improves from 2p$\to$3p.}
\label{tab:all_models}
\begin{tabular}{@{}lcccccr@{}}
\toprule
\textbf{Model} & \textbf{Params} & \textbf{1p} & \textbf{2p} & \textbf{3p} & \textbf{3p H@10} & \textbf{2p$\to$3p} \\
\midrule
\textbf{DELTA-Matched} & 158K & 0.533 & \textbf{0.733} & \textbf{0.738} & \textbf{0.873} & \textbf{+0.005} \\
SBHybrid & 381K & \textbf{0.541} & 0.722 & 0.695 & 0.833 & $-0.027$ \\
GraphGPS & 228K & 0.549 & 0.718 & 0.697 & 0.850 & $-0.021$ \\
DELTA-Full & 293K & 0.524 & 0.711 & 0.692 & 0.836 & $-0.020$ \\
SelfBootstrap & 299K & 0.512 & 0.712 & 0.686 & 0.829 & $-0.026$ \\
GRIT & 197K & 0.460 & 0.712 & 0.644 & 0.783 & $-0.068$ \\
DistMult & 47K & 0.532 & 0.715 & 0.566 & 0.710 & $-0.150$ \\
\bottomrule
\end{tabular}
\end{table}

\paragraph{Standard link prediction (1p).}
\label{app:lp_full}

\begin{table}[h]
\centering
\caption{Standard link prediction (1p) on FB15k-237 top-500 subgraph. Best-validation checkpoint, seed=1.}
\label{tab:lp}
\begin{tabular}{@{}lcccccc@{}}
\toprule
\textbf{Model} & \textbf{Params} & \textbf{MRR} & \textbf{H@1} & \textbf{H@3} & \textbf{H@10} & \textbf{Peak Ep} \\
\midrule
GraphGPS & 228K & \textbf{0.513} & \textbf{0.375} & 0.581 & 0.813 & 200 \\
SBHybrid & 381K & 0.509 & 0.363 & \textbf{0.587} & \textbf{0.816} & 250 \\
DELTA-Matched & 158K & 0.495 & 0.346 & 0.572 & 0.804 & 200 \\
DELTA-Full & 293K & 0.494 & 0.355 & 0.559 & 0.792 & 200 \\
DistMult & 47K & 0.484 & 0.346 & 0.548 & 0.763 & 500+ \\
GRIT & 197K & 0.439 & 0.295 & 0.496 & 0.760 & 200 \\
\bottomrule
\end{tabular}
\end{table}

GraphGPS leads on 1p MRR (0.513 vs.\ DELTA-Matched 0.495, $-0.018$ gap), while DELTA-Matched achieves 97\% of GraphGPS's MRR at 69\% of parameters (158K vs.\ 228K). All subgraph models substantially exceed published full FB15k-237 results (DistMult 0.241, CompGCN 0.355), reflecting the higher link density of the degree-biased subgraph.

%----------------------------------------------------------------------
\section{Full Ablation Results (Phases 67--68)}
\label{app:ablation_full}
%----------------------------------------------------------------------

\begin{table}[h]
\small
\centering
\caption{Phase 67 --- Full FB15k-237 hop-depth ablation (3 seeds $\times$ 200 epochs, top-$k$=128 sparse attention, 1,064K parameters). Values are mean $\pm$ std across all 3 seeds. 2p/3p queries use 10,000 samples each.}
\label{tab:phase67}
\resizebox{\linewidth}{!}{\begin{tabular}{@{}lcccccc@{}}
\toprule
\textbf{Config} & \textbf{LP MRR} & \textbf{H@10} & \textbf{1p MRR} & \textbf{2p MRR} & \textbf{3p MRR} & \textbf{2p$\to$3p$^\dagger$} \\
\midrule
hops=1 & $0.205\pm0.002$ & $0.381\pm0.003$ & $0.257\pm0.004$ & $0.113\pm0.003$ & $0.206\pm0.005$ & $+0.093$ \\
hops=2 & $0.204\pm0.002$ & $0.382\pm0.002$ & $0.252\pm0.003$ & $0.116\pm0.003$ & $0.218\pm0.008$ & $+0.102$ \\
\midrule
$\Delta$ (2$-$1) & $-0.001$ & $+0.001$ & $-0.005$ & $+0.003$ & $\mathbf{+0.012}$ & --- \\
\bottomrule
\end{tabular}}
\vspace{1pt}
{\small $^\dagger$Absolute 2p$\to$3p improvement (3p MRR $-$ 2p MRR). Both conditions show positive improvement from 2p to 3p on the full graph.}
\end{table}

\begin{table}[h]
\small
\centering
\caption{Phase 68 --- Random-subgraph hop-depth ablation (3 seeds $\times$ 500 epochs, $N{=}2{,}242$ uniformly sampled entities, mean degree~7.6, hops=2 adjacency capped at 964K pairs). Values are mean $\pm$ std.}
\label{tab:phase68}
\begin{tabular}{@{}lccccc@{}}
\toprule
\textbf{Config} & \textbf{LP MRR} & \textbf{1p MRR} & \textbf{2p MRR} & \textbf{3p MRR} & \textbf{2p$\to$3p} \\
\midrule
hops=1 & $0.229\pm0.027$ & $0.334\pm0.040$ & $0.242\pm0.057$ & $0.274\pm0.060$ & $+0.032$ \\
hops=2 & $0.233\pm0.019$ & $0.340\pm0.017$ & $0.249\pm0.028$ & $0.284\pm0.031$ & $+0.035$ \\
\midrule
$\Delta$ (2$-$1) & $+0.004$ & $+0.006$ & $+0.007$ & $+0.010$ & --- \\
\bottomrule
\end{tabular}
\end{table}

Per-seed variance for Phase 67 is low: hops=1 seeds span 3p MRR 0.203--0.214; hops=2 seeds 42 and 123 exceed their hops=1 counterparts (3p 0.219--0.227 vs.\ 0.203--0.214), while seed 456 is comparable (0.208). The $+0.003$ 2p gain falls below the $+0.010$ pre-specified threshold, consistent with 2-hop providing incremental benefit primarily at longer depths.

%----------------------------------------------------------------------
\section{Inference Efficiency}
\label{app:inference}
%----------------------------------------------------------------------

\begin{table}[h]
\centering
\caption{Inference timing comparison (mean of 10 runs, CUDA-synchronized).}
\label{tab:timing}
\begin{tabular}{@{}lccc@{}}
\toprule
\textbf{Metric} & \textbf{DELTA-Matched} & \textbf{GraphGPS} & \textbf{Ratio} \\
\midrule
Encoding (once) & 454.1 ms & 8.8 ms & 51.8$\times$ \\
1p per-query & 777.7 $\mu$s & 921.9 $\mu$s & \textbf{0.84$\times$} \\
2p per-query & 1,380 $\mu$s & 1,476 $\mu$s & \textbf{0.94$\times$} \\
3p per-query & 1,251 $\mu$s & 1,371 $\mu$s & \textbf{0.91$\times$} \\
Training (total) & 3,782 s & 110 s & 34.2$\times$ \\
\bottomrule
\end{tabular}
\end{table}

Encoding---a one-time step including 2-hop edge adjacency construction---is 51.8$\times$ slower than GraphGPS. However, \emph{per-query scoring} is 0.84--0.94$\times$ GraphGPS: DELTA is faster per query. In deployment scenarios where the graph is encoded once and queries are scored many times, DELTA's amortized cost is comparable to or better than GraphGPS. Training cost (34.2$\times$) is a one-time expense dominated by edge adjacency computation.

%----------------------------------------------------------------------
\section{Scaling Analysis}
\label{app:scaling}
%----------------------------------------------------------------------

\begin{table}[h]
\centering
\caption{Scaling analysis: DELTA (1-layer) vs.\ DistMult across entity counts. All models on FB15k-237 subgraphs, seed=42. The full FB15k-237 row (Phase 67) uses top-$k$=128 sparse attention and 1,064K parameters; DM result is from published literature.}
\label{tab:scaling}
\begin{tabular}{@{}cccccccc@{}}
\toprule
\textbf{$N$} & \textbf{Triples} & \textbf{$|E_{\text{adj}}|$} & \multicolumn{2}{c}{\textbf{Test MRR}} & \textbf{$\Delta$} & \textbf{DELTA time} & \textbf{DM time} \\
\cmidrule(lr){4-5}
& & & DELTA & DM & & & \\
\midrule
500 & 9,700 & 256K & 0.495 & 0.484 & +0.011 & 3,782s & 110s \\
2,000 & 62,700 & 15.2M & 0.334 & 0.319 & +0.015 & 5,788s & 72s \\
5,000 & 152,800 & 63.0M & 0.247 & 0.224 & +0.023 & 21,688s & 1,784s \\
14,541 & 310,116 & 29.8M$^\dagger$ & 0.205 & 0.241$^*$ & n/a$^*$ & --- & --- \\
\bottomrule
\end{tabular}
\vspace{1pt}
{\small $^\dagger$hops=1 adjacency. $^*$\textbf{Not directly comparable:} the published DistMult result~\cite{yang2015embedding} uses tuned hyperparameters and extended training time; our DELTA result at $N{=}14{,}541$ uses a fixed training budget. The $\Delta{=}{-0.036}$ entry does \emph{not} establish that DELTA underperforms DistMult under matched conditions at full scale.}
\end{table}

\paragraph{Genuine but modest advantage at scale.}
DELTA consistently outperforms DistMult at all controlled scales tested ($N{=}500$ to $N{=}5{,}000$), with the gap growing from $+0.011$ to $+0.023$. A fair comparison at $N{=}14{,}541$ requires matched training conditions and is left to future work.

\paragraph{Depth is the scaling bottleneck.}
At $N{=}2{,}000$, 3-layer DELTA catastrophically over-smooths (MRR $0.002$, near random), while 1-layer DELTA surpasses DistMult. Residual gating eliminates the failure mode but gates freeze at ${\sim}10\%$ layer / $90\%$ residual, confirming that additional depth provides no benefit at this scale.

\paragraph{Top-$k$ sparse attention.}
At $k{=}128$, sparse attention on 63M adjacency pairs matches full-softmax quality exactly (MRR $0.2472$ vs.\ $0.2457$), while $k{=}64$ degrades by $-5.6\%$. Subsampling to ${\sim}48\%$ retention at $N{=}5{,}000$ slightly improves MRR ($+0.007$), suggesting that a moderate number of high-quality neighbors outperforms access to all neighbors.

%----------------------------------------------------------------------
\section{Attention Temperature Asymmetry}
\label{app:temperature}
%----------------------------------------------------------------------

We introduce learnable per-head temperature parameters $\tau^{(h)} = \exp(\theta^{(h)})$ that scale attention scores before softmax (Eq.~1 in the main paper). Higher $\tau$ produces sharper (more selective) attention; lower $\tau$ produces softer (more uniform) attention. We report these as observational findings; causal attribution requires controlled interventions (temperature freezing, per-head ablations) that we leave to future work.

\paragraph{Edge--node temperature divergence.}
Across all 20+ configurations tested, \textbf{edge attention temperatures drift upward} (from 4.0 to 4.4--5.0) while \textbf{node attention temperatures drift downward} (from 4.0 to 3.5--3.7). The model autonomously discovers that edges benefit from selective, discriminative attention while nodes benefit from broader, averaging attention.

\paragraph{Layer 0 is always dead.}
Regardless of initialization, first-layer attention heads remain near-uniform (entropy $>0.95$). Initial entity embeddings carry no relational information worth selectively attending to; useful attention emerges only after the first reconciliation bridge provides cross-stream context. This suggests the bridge's sequential cascade may initiate representational differentiation that later layers exploit---a causal question left for future bridge-ablation experiments.

\paragraph{Temperature controls operating mode.}
Three distinct operating regimes emerge:
\begin{enumerate}
    \item \textbf{LP-optimized} (asymmetric: node$=$2, edge$=$6): LP MRR 0.491, sacrifices multi-hop.
    \item \textbf{Balanced} (uniform temp$=$4.0): Best 3p MRR (0.402), moderate LP.
    \item \textbf{Annealed} (node temp $4.0 \to 2.0$ during training): Breaks the 3p ceiling (0.415).
\end{enumerate}

Temperature annealing creates compositional representations in early epochs that static initialization cannot replicate---training \emph{trajectory} matters, not just final temperature. Multi-seed validation (3 seeds) confirms LP improvements are statistically robust (MRR $0.484 \pm 0.005$); multi-hop temperature effects showed high cross-seed variance that resolved with larger query counts (10,000 samples, reducing std by 66--84\%).

%----------------------------------------------------------------------
\section{Detailed Limitations}
\label{app:limitations}
%----------------------------------------------------------------------

\textbf{Scale and benchmark coverage (critical).}
Compositional reasoning experiments use a 500-entity degree-biased subgraph (mean degree $\approx 19.7$), far denser than the full graph ($\approx 4.1$). The strong multi-hop results (0.758--0.790 MRR, 5-hop) are demonstrated on this small subgraph. On the complete 14,541-entity FB15k-237 graph, DELTA achieves 1p MRR $\approx 0.205$, substantially below published SOTA: RotatE 0.338, CompGCN 0.355, NBFNet 0.415. Evaluation on other standard benchmarks (YAGO3-10, WN18RR~\cite{dettmers2018convolutional}, ogbl-wikikg2) is an important open direction.

\textbf{Baseline coverage (important).}
We compare primarily against general graph transformers (GraphGPS, GRIT) and DistMult. Specialized path-based KG reasoners (NBFNet~\cite{zhu2021neural}, A*Net~\cite{zhu2022astar}) are absent. Our multi-hop evaluation uses internally generated chain queries with soft entity traversal---not the established BetaE/ConE/GQE complex-query benchmarks, which use different query generation protocols and logical operators. Results are not directly comparable to published multi-hop reasoning numbers from those systems.

\textbf{Full-scale multi-hop evaluation (important).}
The Phase 67 full-graph ablation evaluates hops=1 vs.\ hops=2 but does not include the full 2-hop$\to$5-hop depth sweep from the subgraph experiments. Whether DELTA's depth-monotonic improvement persists on the full sparse graph has not been directly validated.

\textbf{Computational cost (notable).}
The one-time encoding overhead (51.8$\times$ GraphGPS) limits applicability to static graphs. The full-graph (Phase 67) experiment required approximately 3 days on an NVIDIA H200 GPU with 141 GB VRAM to evaluate 6 conditions (3 seeds $\times$ 2 hop settings $\times$ 200 epochs), with seeds running in parallel.

\textbf{Causal claims (methodological).}
The attention temperature asymmetry is a consistent observational pattern, not a causally established property. Causal attribution requires controlled interventions (temperature freezing, head-type ablations) not performed here. BrainEncoder (Appendix~\ref{app:brain}) and domain transfer (Appendix~\ref{app:transfer}) results are preliminary signals from small experiments and should not be treated as validated findings.

\textbf{Subgraph selection bias (methodological).}
Using top-degree entities for primary experiments biases toward entities with rich relational context and a highly saturated edge adjacency graph, which may affect which architectural choices appear beneficial relative to sparser realistic settings.

%----------------------------------------------------------------------
% NeurIPS Paper Checklist
%----------------------------------------------------------------------
\newpage
\section*{NeurIPS Paper Checklist}

\begin{enumerate}

\item {\bf Claims}
\item[] Question: Do the main claims made in the abstract and introduction accurately reflect the paper's contributions and scope?
\item[] Answer: \answerYes{}
\item[] Justification: The abstract and introduction claim (1) that edge-to-edge attention over multi-hop adjacency improves compositional reasoning on knowledge graphs, and (2) that DELTA demonstrates consistent improvements on synthetic, full-graph, and subgraph benchmarks. All claims are directly supported by the experimental results in \S\ref{sec:experiments} with numerical evidence. Scope caveats (SOTA LP gap, dense-subgraph null ablation) are stated explicitly in the abstract and \S\ref{sec:discussion}.

\item {\bf Limitations}
\item[] Question: Does the paper discuss the limitations of the work?
\item[] Answer: \answerYes{}
\item[] Justification: Limitations are addressed in \S\ref{sec:discussion} (summary paragraph) and Appendix~\ref{app:limitations} (detailed): the gap between SOTA and DELTA's full-graph LP; the Phase 66 dense-subgraph ablation finding all conditions statistically indistinguishable; the preliminary nature of BrainEncoder and domain transfer results; and missing baselines (NBFNet, A*Net).

\item {\bf Theory Assumptions and Proofs}
\item[] Question: For each theoretical result, does the paper provide complete proofs of all formal claims?
\item[] Answer: \answerNA{}
\item[] Justification: The paper does not make formal theoretical claims requiring proofs. The $O(E^{0.97})$ scaling claim is validated empirically (\S\ref{sec:edge_adj}). Architectural design choices are motivated and validated experimentally.

\item {\bf Experimental Result Reproducibility}
\item[] Question: Does the paper fully describe the experimental setup to enable re-implementation?
\item[] Answer: \answerYes{}
\item[] Justification: \S\ref{sec:exp_setup} specifies dataset (FB15k-237, publicly available), subgraph construction (top-$k$ degree filter), model hyperparameters (hidden dimensions, number of layers, hop depth, $\text{lr}{=}0.003$, batch size 4,096, label smoothing 0.1, DropEdge rate), evaluation metrics, and number of seeds per experiment. Phase identifiers map to distinct experimental conditions described in the text. Code covering edge adjacency construction, top-$k$ sparse attention, and multi-hop query generation will be released upon acceptance.

\item {\bf Open Access to Data and Code}
\item[] Question: Does the paper provide open access to the data and code, with instructions for reproducibility?
\item[] Answer: \answerNo{}
\item[] Justification: Code will be released upon acceptance. All datasets used (FB15k-237, WN18RR) are publicly available benchmarks. The paper provides sufficient architectural and hyperparameter detail to enable re-implementation without the code.

\item {\bf Experimental Setting/Details}
\item[] Question: Does the paper specify all the training and test details (e.g., data splits, hyperparameters, how they were chosen)?
\item[] Answer: \answerYes{}
\item[] Justification: \S\ref{sec:exp_setup} specifies optimizer (Adam, $\text{lr}{=}0.003$), batch size (4,096), loss (1-vs-all BCE with label smoothing 0.1), early stopping (patience 10), evaluation metrics (filtered MRR, Hits@$k$), number of epochs per phase (200--500 depending on experiment), and seeds used. Hyperparameters were selected based on prior GNN work and kept fixed across all models for fair comparison.

\item {\bf Experiment Statistical Significance}
\item[] Question: Does the paper report error bars suitably and correctly defined or other appropriate information about the statistical significance of the experiments?
\item[] Answer: \answerYes{}
\item[] Justification: Multi-seed experiments (Phases 34, 66, 67, 68) report mean $\pm$ std across 3--5 seeds throughout Tables~\ref{tab:synthetic}, \ref{tab:ablation_hop}, \ref{tab:multiseed}, \ref{tab:cross_density}, \ref{tab:phase67}, and \ref{tab:phase68}. Error bars represent standard deviation across seeds. The Phase 66 null result and the Phase 67 threshold comparison explicitly use standard deviation overlap to assess statistical significance.

\item {\bf Experiments Compute Resource}
\item[] Question: For each experiment, does the paper provide sufficient information on the computer resources (type of compute, memory, time) needed to reproduce the experiments?
\item[] Answer: \answerYes{}
\item[] Justification: Synthetic benchmarks (Phase 34) ran on an H100 GPU (mentioned in abstract and \S\ref{sec:exp_synthetic}). Subgraph experiments (Phases 66, 68) used 500 epochs; full-graph experiments (Phase 67) used 200 epochs with top-$k{=}128$ sparse attention. Appendix~\ref{app:inference} reports training time: DELTA-Matched total training is 3,782s vs.\ GraphGPS 110s. Per-query inference is 0.84--0.94$\times$ GraphGPS after the one-time 454ms encoding step.

\item {\bf Code Of Ethics}
\item[] Question: Does the research conducted in the paper conform, in every way, with the NeurIPS Code of Ethics?
\item[] Answer: \answerYes{}
\item[] Justification: This work involves publicly available knowledge graph benchmarks (FB15k-237, WN18RR) with no human subjects, no sensitive data, and no potential for direct harm. No crowdsourcing was used.

\item {\bf Broader Impacts}
\item[] Question: Does the paper discuss both potential positive societal impacts and negative societal impacts of the work?
\item[] Answer: \answerNA{}
\item[] Justification: This is foundational research on graph neural network architectures for knowledge graph reasoning. There are no foreseeable direct negative societal impacts beyond general concerns about AI misuse not specific to this contribution.

\item {\bf Safeguards}
\item[] Question: Does the paper describe safeguards that have been put in place for responsible release of data or models?
\item[] Answer: \answerNA{}
\item[] Justification: All datasets used are already publicly available standard benchmarks. No new datasets or high-risk models are released with this submission.

\item {\bf Licenses for Existing Assets}
\item[] Question: Are the creators or original owners of assets (e.g., code, data, models) used in the paper properly credited?
\item[] Answer: \answerYes{}
\item[] Justification: FB15k-237~\cite{toutanova2015representing} is cited. GraphGPS~\cite{rampavsek2022recipe} and GRIT~\cite{ma2023graph} are cited as baselines. DistMult~\cite{yang2015embedding}, RotatE~\cite{sun2019rotate}, NBFNet~\cite{zhu2021neural}, and A*Net~\cite{zhu2022astar} are properly attributed. All referenced architectures and datasets have corresponding bibliography entries.

\item {\bf New Assets}
\item[] Question: Are new assets introduced in the paper documented with appropriate information?
\item[] Answer: \answerNA{}
\item[] Justification: No new datasets or pre-trained models are released with this submission. Code will be released upon acceptance.

\item {\bf Crowdsourcing and Research with Human Subjects}
\item[] Question: For crowdsourcing experiments and research with human subjects, does the paper include the necessary information (e.g., description of tasks, payments, and any instructions given to participants)?
\item[] Answer: \answerNA{}
\item[] Justification: No crowdsourcing or human subject experiments were conducted.

\item {\bf Institutional Review Board (IRB) Approvals}
\item[] Question: Does the paper describe whether IRB approvals (or equivalent) were obtained?
\item[] Answer: \answerNA{}
\item[] Justification: No human subjects research was conducted.

\item {\bf Declaration of LLM Usage}
\item[] Question: Does the paper describe the usage of LLMs if they are an important, original, or non-standard component of the core methods?
\item[] Answer: \answerNA{}
\item[] Justification: LLMs were not used as any component of the core research methods. LLM-based tools were used for writing assistance (editing, grammar, drafting) and code assistance during development, but do not constitute an original or non-standard part of the methodology.

\end{enumerate}

\end{document}
